%%=============================================================================
%% Technologische context en keuze
%%=============================================================================

\chapter{\IfLanguageName{dutch}{Technologie en verantwoording}{Technology and rationale}}%
\label{ch:technologie}

%% Richtlengte: 2 tot 3 bladzijden.
%%
%% LET OP: DIT IS GEEN LITERATUURSTUDIE.
%% Je hoeft geen overzicht te schrijven van wat er in het vakgebied bestaat,
%% en je hoeft geen wetenschappelijke bronnen te verzamelen. De vraag hier is
%% praktisch: welke technologie heb je gebruikt, welke alternatieven waren er,
%% en waarom is het deze geworden?
%%
%% Dit is ook wat de jury je tijdens je verdediging zal vragen. Wie het hier
%% opschrijft, staat daar sterker.
%%
%% WEL: vermeld de bronnen die je echt gebruikt hebt. Voor een project zijn dat
%% meestal officiële documentatie, handleidingen, blogberichten van makers,
%% normen, of een boek. Vermeld ze omdat je lezer ze moet kunnen terugvinden,
%% niet omdat een lijst nu eenmaal lang moet zijn. Zie de uitleg onderaan dit
%% bestand en de voorbeelden in bronnen.bib.

Vervang deze tekst door een korte inleiding op het hoofdstuk: welke
technologische omgeving bestond er al bij het bedrijf, en welke ruimte had je
om te kiezen?

\section{\IfLanguageName{dutch}{Bestaande omgeving}{Existing environment}}%
\label{sec:bestaande-omgeving}

%% Waar stap je in? Welke talen, frameworks, databanken, diensten en
%% werkwijzen waren er al? Wat lag vast en wat mocht je zelf beslissen?

\section{\IfLanguageName{dutch}{Vergelijking van de mogelijkheden}{Comparison of the options}}%
\label{sec:vergelijking}

%% Vergelijk twee tot vier reële mogelijkheden op criteria die er voor dit
%% project toe doen: aansluiting bij de bestaande omgeving, leercurve,
%% ondersteuning, licentie, prestaties, onderhoudbaarheid, kostprijs.
%%
%% Een tabel werkt hier goed. Verzin geen criteria om een tabel te vullen:
%% drie criteria die echt gewogen hebben, zijn beter dan acht die je verzon.

\begin{table}[htbp]
  \centering
  \begin{tabular}{lccc}
    \toprule
    \textbf{Criterium} & \textbf{Optie A} & \textbf{Optie B} & \textbf{Optie C} \\
    \midrule
    Aansluiting bij de bestaande omgeving & ++ & +  & -- \\
    Leercurve binnen tien dagen           & +  & ++ & -- \\
    Ondersteuning en documentatie         & ++ & +  & +  \\
    \bottomrule
  \end{tabular}
  \caption[Vergelijking van de overwogen technologieën]{\label{tab:vergelijking}%
    Vergelijking van de overwogen technologieën. Vervang de criteria en de
    beoordeling door je eigen afweging, en licht ze in de tekst toe.}
\end{table}

\section{\IfLanguageName{dutch}{Vergelijking met de opleiding}{Comparison with the programme}}%
\label{sec:vergelijking-opleiding}

%% Verwacht onderdeel van je verdediging: hoe verhoudt de technologie die je
%% hier gebruikte zich tot wat je in de opleiding gezien hebt? Wat was
%% hetzelfde, wat was nieuw, en wat bleek in de praktijk anders te werken dan
%% in de les?

\section{\IfLanguageName{dutch}{Gemaakte keuze}{The choice made}}%
\label{sec:keuze}

%% Eén heldere paragraaf: dit is het geworden, en dit is de reden. Vermeld
%% ook eerlijk wat je ervoor hebt moeten opgeven.
%%
%% Verwijs naar je bronnen met \autocite{sleutel} of \textcite{sleutel}. De
%% sleutels staan in bronnen.bib. Zet de verwijzing binnen de zin, vlak bij de
%% bewering waar ze bij hoort - niet achteraan de paragraaf. Hieronder een
%% voorbeeld van beide vormen; vervang ze door je eigen bronnen.

Bij het beoordelen van leesbaarheid en onderhoudbaarheid van de code werd
uitgegaan van de gangbare vuistregels uit het vakgebied~\autocite{Martin2008}.
De officiële documentatie van het gekozen framework~\autocite{DockerDocs2026}
diende als voornaamste naslagwerk. \textcite{Fowler2018} gebruik je wanneer de
naam van de auteur zelf een rol speelt in je zin.

%%---------- Bronnen vermelden in een projectrapport ---------------------------
%%
%% Welke bronnen horen in je bronnenlijst?
%%   - officiële documentatie van de talen, frameworks en diensten die je gebruikt
%%   - handleidingen, blogberichten of video's waarop je je aanpak gebaseerd hebt
%%   - normen en standaarden waaraan je je gehouden hebt
%%   - een boek of artikel, als je er een gebruikt hebt
%%   - interne documenten van je bedrijf: vermeld ze, maar neem geen
%%     vertrouwelijke inhoud over zonder toestemming van je mentor
%%
%% Wat er NIET in hoort:
%%   - een gesprek met een AI-hulpmiddel. Dat is geen bron die iemand kan
%%     nagaan; het hoort thuis in de bijlage over je AI-gebruik.
%%   - bronnen die je niet gelezen hebt. Een lange lijst maakt geen indruk;
%%     een lijst met werken die nergens aangehaald worden, wel - de verkeerde.
%%
%% Vermeld bij een webbron ALTIJD de datum waarop je ze geraadpleegd hebt
%% (het veld urldate in bronnen.bib): documentatie verandert.
